% 热木星（综述）
% license CCBYSA3
% type Wiki

本文根据 CC-BY-SA 协议转载翻译自维基百科\href{https://en.wikipedia.org/wiki/Hot_Jupiter}{相关文章}。

\begin{figure}[ht]
\centering
\includegraphics[width=6cm]{./figures/d9160f8109011126.png}
\caption{一幅艺术家绘制的近距离绕其母星运行的热木星示意图} \label{fig_HoJu_1}
\end{figure}
热木星（有时也称为热土星）是一类气态巨行星系外行星，被认为在物理性质上与木星相似（即木星类行星），但其轨道周期极短（(P<10) 天）$^{[1]}$。由于它们与母星距离极近、表面与大气温度很高，因而获得了“热木星”这一非正式名称$^{[2]}$。

在所有系外行星中，热木星是最容易通过径向速度法探测到的类型之一，这是因为它们在母星运动中引起的摆动幅度相对较大、变化也更快，明显强于其他已知类型的行星。其中最著名的热木星之一是 51 Pegasi b。它于 1995 年被发现，是第一颗被证实绕类太阳恒星运行的系外行星。51 Pegasi b 的轨道周期约为 4 天$^{[3]}$。
\subsection{一般特征}
\begin{figure}[ht]
\centering
\includegraphics[width=6cm]{./figures/e311dba254552faa.png}
\caption{截至 2014 年 1 月 2 日已发现的热木星（分布在左侧边缘，其中包括大多数通过凌日法探测到的行星，以黑点表示）} \label{fig_HoJu_2}
\end{figure}
\begin{figure}[ht]
\centering
\includegraphics[width=6cm]{./figures/c9637a190eec93c9.png}
\caption{具有隐藏水分子的热木星$^{[4]+}$} \label{fig_HoJu_3}
\end{figure}